%% kapa: exam=1 round=33 year=2022 period=2 form=A subject=appraisal_law subject_ko=감정평가관계법규 seq=11 number=11
%% source: Q-Net 감정평가사 기출문제 (cst00309) artlSeq=5209997 file=raw/1008/5209997/2022년도 제33회 감정평가사 1차시험 2교시(감정평가 관계법규, 회계학).pdf
%% conversion: kordoc 4.24.51 -> tools/gen1.py (auto); review=proofread-sonnet; flags=none
\begin{question}{11}
국토의 계획 및 이용에 관한 법령상 광역도시계획의 수립에 관한 설명으로 옳지 않은 것은?
\begin{choices}
  \item 국토교통부장관은 시·도지사가 요청하는 경우 관할 시·도지사와 공동으로 광역도시계획을 수립할 수 있다.
  \item 시·도지사가 광역도시계획을 수립하는 경우 미리 공청회를 열어 주민과 관계 전문가 등으로부터 의견을 들어야 한다.
  \item 국토교통부장관은 관계 행정기관의 장에게 광역도시계획의 수립을 위한 기초조사에 필요한 자료를 제출하도록 요청할 수 있다.
  \item 시·도지사가 광역도시계획을 수립하는 경우 미리 관계 중앙행정기관과 협의한 후 중앙도시계획위원회의 심의를 거쳐야 한다.
  \item 시·도지사가 광역도시계획의 승인을 받으려는 때에는 광역도시계획안에 기초조사 결과를 포함한 서류를 첨부하여 국토교통부장관에게 제출해야 한다.
\end{choices}
\end{question}
\answer{4}
